% Generated draft: requires mathematical and source review.
\documentclass[11pt]{article}
\usepackage[margin=25mm]{geometry}
\usepackage{fontspec}
\setmainfont{Latin Modern Roman}
\usepackage{amsmath,amssymb,mathtools,unicode-math}
\setmathfont{Latin Modern Math}
\usepackage{enumitem,needspace,xurl,hyperref}
\setlength{\parindent}{0pt}
\setlength{\parskip}{5pt}
\setlength{\emergencystretch}{3em}
\newcommand{\N}{\mathbb{N}}
\newcommand{\Z}{\mathbb{Z}}
\newcommand{\Q}{\mathbb{Q}}
\newcommand{\R}{\mathbb{R}}
\newcommand{\C}{\mathbb{C}}
\newcommand{\F}{\mathbb{F}}
\newcommand{\A}{\mathbb{A}}
\newcommand{\PP}{\mathbb{P}}
\newcommand{\E}{\mathbb{E}}
\newcommand{\eps}{\varepsilon}
\newcommand{\dd}{\,\mathrm{d}}
\begin{document}
\section*{UnsolvedMath ID: 701030}
\textbf{Source problem code:} catalogue entry
\subsection{Definitions and mathematical statement}
A group \(G\) is linearly right-orderable if there exists a strict total order \(<\) on \(G\) such that
\[
\forall x,y,z\in G,\qquad x<y\Longrightarrow xz<yz.
\]
A subgroup \(H\leq G\) is cyclic if \(H=\langle h\rangle\) for some \(h\in H\). Here \(H<G\) means that \(H\) is a proper subgroup, including the trivial subgroup. A group is non-abelian simple if it is non-abelian and its only normal subgroups are \(\{1\}\) and itself.

Interpreting the source's parenthetical words ``non-abelian simple'' as requirements, the question is whether there exists a non-abelian simple group \(G\) admitting a right-invariant strict total order and satisfying
\[
\forall H\leq G,\qquad H\neq G\Longrightarrow \exists h\in H\;\bigl(H=\langle h\rangle\bigr).
\]
The source places ``non-abelian simple'' in parentheses without explaining their logical role. Thus the displayed formulation explicitly adopts the reading that these are required properties; whether the source instead intends a broader existence question with a parenthetical qualification remains unresolved. No cardinality assumption is supplied.
\subsection{Short English statement}
Does there exist a non-abelian simple group with a right-invariant total order in which every proper subgroup is cyclic? This reads the source's parenthetical simplicity and noncommutativity conditions as requirements; their intended logical role requires review.
\subsection{Sources}
\begin{itemize}
\item Supplied local UnsolvedMath registry record, permanent ID \(701030\), exact source problem code \texttt{catalogue entry}, titled \emph{Kourovka Notebook Problem 14.45}. The registry labels the problem ``open''; this is supplied source metadata, not an independently verified status.
\item \emph{The Kourovka Notebook (Unsolved Problems in Group Theory)}, 14th issue (1999), Problem 14.45, proposed by U. E. Kaljulaid, as attributed in the supplied record.
\item The supplied background also identifies the 21st issue edition (2026), \texttt{arXiv:1401.0300v48}. This edition information is reproduced from the record and has not been independently verified.
\item The supplied collection description credits E. I. Khukhro and V. D. Mazurov as editors of the Kourovka Notebook collection covering Issues 1--21 (1965--2026).
\end{itemize}
\end{document}
